\section{Limited Client Resources}

% TODO: use nas-bench 101 blog post as citation for explanation on why NAS is so expensive
NAS is inherently computationally demanding, because a search may train and compare many candidate architectures. Even in a centralised setting it is difficult to do NAS.

NAS-Bench-101 made repeated NAS experiments practical through table look-ups by precomputing more than five million trained models for 423,000 unique convolutional architectures, a dataset that required over 100 TPU-years of computation~\cite{ying2019nasbench101reproducibleneuralarchitecture}. This scale illustrates the cost that candidate evaluation can impose before FL-specific constraints are considered.

FL compounds this cost because cross-device clients vary in compute performance, memory, bandwidth and availability, and these capabilities are typically lower than those of data-centre nodes~\cite{fl_advances_and_open_problems_2021,fl_in_practice_reflections_2024}. This thematic analysis identified four resource-related sub-themes: limited computation and memory, limited networking, energy consumption and client unreliability.

These constraints are generally less severe in cross-silo FL because its clients are usually fewer, more powerful, more reliable and addressable. This eliminates many cross-device system challenges, although computation or communication can still be the primary bottleneck~\cite{fl_advances_and_open_problems_2021}.

\subsection{Limited Computational Power} % challenges and solutions related to computation

\subsubsection{Challenge 1.1: Evaluating Candidate Architectures on Constrained Computational Resources}

NAS-in-FL can become infeasible when the model that a client must evaluate in one communication round exceeds its available memory capacity or the evaluation cannot be completed within the time available for a communication round. This concern is central to approaches that target mobile and edge clients~\cite{fedoras_2022,fednas_over_hetero_mobile_devices_2022,dfnas_2020} where client memory and computational power is rather limited. The problem is particularly acute for Differentiable Supernet-Based strategies that distribute a shared supernet and require clients to update multiple mixed candidate operations locally~\cite{fednas_2021,dp-fnas_2020,feathers_2023,fed_auto_mri_2023}. Other search strategy types usually assign one discrete or generated candidate at a time, but it remains challenging to evaluate even just one candidate model in a feasible amount of time.

\noindent\textbf{Solution 1.1a: Build the search space from inexpensive components.} Practitioners can apply this solution to every search strategy type by restricting candidate architectures to inexpensive operations, blocks or modules, depending on the search space type. Prominent examples include depthwise-separable convolutions and mobile-oriented CNN blocks~\cite{fed_auto_mri_2023,intrusion_detection_fednas_2024}. This makes client-side evaluation cheaper, but may exclude expensive components that increase prediction accuracy.

\noindent\textbf{Solution 1.1b: Train only a subnet per client at a time.} Supernet-based search strategies can be modified so that clients train and evaluate a subnet corresponding to a subset of candidate architectures instead of a weighted mixture of all candidate operations, reducing the computational burden on individual clients~\cite{dfnas_2020}. The subnet size can be adjusted dynamically to the computational budget of each client, with the most extreme case being evaluation of a single path~\cite{fedoras_2022,divide-and-conquer_fednas_2023,optimize_fednas_edge_2021,efnas_2024,load_monitoring_fednas_2025}. However, sampling only part of the supernet in each round can evaluate candidate operations unevenly and thereby bias comparisons between candidate architectures~\cite{fedoras_2022,hapfnas_2025}.

\noindent\textbf{Solution 1.1c: Search trainable adaptation modules over a frozen backbone.} When a suitable pretrained model is available, clients can keep its backbone fixed and search a sparse combination of lightweight adaptation modules. Only the selected modules and their architecture variables require local optimisation and exchange, reducing both client computation and update size~\cite{fednaspet_2023}. The benefit depends on how well the fixed backbone transfers to the clients' tasks.

\subsubsection{Challenge 1.2: Reducing the Cost of the Federated Search Across Rounds}

Even when evaluating candidates on individual clients is feasible, keeping the overall duration of the federated architecture search within a practical time frame can be difficult because the search must compare architectural alternatives across repeated rounds in addition to optimising model weights. Black-box strategies may train several discrete candidates through separate federated evaluations~\cite{decnas_2022,rt_fed_evo_nas_2020}, while weight-sharing strategies repeatedly update an over-parameterised supernet or its sampled subnets~\cite{fednas_2021,superfednas_2024}. A separate post-search training cost applies only when the selected architecture cannot reuse search-time weights, as discussed in Challenge 1.3.

\noindent\textbf{Solution 1.2a: Dynamic candidate evaluation budgets.} This solution applies to black-box search strategies that compare discrete candidates through independent training, including greedy pruning, evolutionary and genetic algorithms, reinforcement learning and Bayesian optimisation. Similar to multi-armed-bandit AutoML approaches ~\cite{automl_book_hpo_2019}, each candidate initially receives only a small budget of federated rounds, after which weak candidates can be discarded and promising candidates trained for longer.

For example, \cite{decnas_2022,fednas_over_hetero_mobile_devices_2022} observes that the best candidate usually becomes distinguishable within the first couple of communication rounds, drops the remaining candidates and increases the round budget in later search iterations. Evolutionary searches similarly use a fixed small number of rounds or early stopping to limit candidate evaluation~\cite{multi_objective_evo_fl_2020,gpt4-boosted_pso_fednas_2024}.

\noindent\textbf{Solution 1.2b: Prune the supernet progressively.} Supernet-based strategies can periodically remove weak candidate operations from the shared supernet, so fewer alternatives remain stored and trained in subsequent federated rounds~\cite{dfnas_2020,peaches_2024}. Unlike dynamic candidate evaluation budgets, this changes the search space itself rather than the number of rounds allocated to independently trained candidates. Its savings increase as the search progresses, but the initial rounds still incur the cost of the complete supernet and early pruning can permanently remove useful operations.

\noindent\textbf{Solution 1.2c: Predict candidate performance.} Performance prediction is a standard AutoML technique for reducing the number of costly configuration evaluations~\cite{automl_book_hpo_2019}. In NAS-in-FL, candidate-based strategies can evaluate a small set of architectures and train a predictor to estimate the quality of the remaining candidates~\cite{hapfnas_2025,energy_fednas_2025}. Additional care is required because private and non-IID client data can make architecture rankings client-specific. Predictors can therefore be trained separately for each client~\cite{hapfnas_2025}, or learned collaboratively by aggregating their parameters without transmitting the locally evaluated architectures~\cite{federated_bayesian_optimization_nas_2023}. This reduces the number of candidates evaluated by clients, but its reliability depends on whether the evaluated architectures represent the relevant parts of the search space.

\noindent\textbf{Solution 1.2d: Evaluate candidates in parallel.} Since black-box search strategies generate several distinct candidates per iteration, the candidates can be trained and evaluated concurrently by separate client groups~\cite{decnas_2022,rt_fed_evo_nas_2020,energy_fednas_2025}.

Supernet-based strategies can instead sample several subnets, allocate each to a client cluster and merge their updates back into the shared supernet~\cite{finch_2024}. This potentially reduces elapsed search time when enough clients are available, but comparisons can become biased when the groups evaluate candidates on different data distributions~\cite{decnas_2022}.

\noindent\textbf{Solution 1.2e: Reuse supernet weights for candidate evaluation.} Weight-sharing strategies can train common operation weights and allow each candidate to inherit the subset associated with its architecture~\cite{fl-agnns_2023,optimize_fednas_edge_2021}. This avoids training every candidate independently from scratch, although inherited-weight performance can rank candidates differently from full training. A short incremental federated fine-tuning stage can provide a higher-fidelity estimate for a smaller set of promising candidates~\cite{medical_image_segmentation_fednas_2024}.

\noindent\textbf{Solution 1.2f: Supervise concurrent search agents.} Several search agents can run in parallel while a supervisor compares their relative losses and adjusts the momentum applied by each agent. The final architecture can then be selected according to prediction error, training time and model size~\cite{supervised_fednas_2023}. This can smooth and accelerate convergence when parallel resources are available, but requires maintaining several search trajectories.

\subsubsection{Challenge 1.3: Post-Search Retraining Duplicates Federated Computation}

Some NAS strategies use search-time training only to select an architecture and then train the selected model from a new initialisation. In NAS-in-FL, this adds another sequence of client-side epochs and aggregation rounds after clients have already spent computation and communication on the search~\cite{medical_image_segmentation_fednas_2024,privacy_preserving_fednas_for_edge_2025}. This challenge applies whenever the search produces an architecture without reusable final weights and therefore delays the availability of a deployable model.

\noindent\textbf{Solution 1.3a: Reuse weights learned during search.} Weight-sharing supernet strategies can initialise the selected subnet with the corresponding operation weights learned during search and perform only limited final fine-tuning~\cite{medical_image_segmentation_fednas_2024}. This reuses client computation that has already been spent on the selected operations, although weights learned jointly with competing subnets may be a poor initialisation for the final model.

\noindent\textbf{Solution 1.3b: Produce a deployable model during search.} One-stage supernet-based strategies can jointly optimise architecture selection and model weights while concentrating the architecture distribution or pruning weak operations, so the search terminates with a discrete model and trained weights~\cite{dfnas_2020,fed_meta_nas_2025}. Reinforcement-learning strategies over a weight-sharing supernet can similarly turn later search rounds into fine-tuning as their sampling probabilities concentrate~\cite{perfedrlnas_2024}. This eliminates the separate retraining stage, but the final path may remain undertrained if it received too few updates while sharing weights with competing paths.

\subsubsection{Challenge 1.4: Separate Searches for Deployment Targets Multiply Cost}

A NAS-in-FL deployment may require different architectures for several latency or computation budgets. Searching and training independently for each target repeats the federated procedure, so the required client computation and communication grow with the number of device classes~\cite{superfednas_2024}.

\noindent\textbf{Solution 1.4a: Reuse one supernet across deployment targets.} A weight-sharing supernet containing subnets of different depths and widths can be trained once through FL. Clients can then use predictor-guided local search to select subnets for their deployment constraints without further federated training. This amortises the training cost across targets and has reduced it elevenfold for 20 targets~\cite{superfednas_2024}. Co-training many subnets can, however, cause weight-sharing interference and leave rarely sampled parameters stale~\cite{superfednas_2024}.

\subsubsection{Challenge 1.5: Split-Model Search Leaves Client Computation Idle}

In split-model training, different clients or device groups execute consecutive model segments and exchange intermediate activations and gradients. The resulting sequential dependency leaves clients waiting while other segments execute. NAS-in-FL methods that distribute a search model in this manner inherit these idle periods during architecture search~\cite{nasfly_2024}.

\noindent\textbf{Solution 1.5a: Search locally during split-training idle periods.} Auxiliary input and output components can turn an assigned model segment into a locally executable supernet. During otherwise idle periods, the client samples and trains alternative branches, using the corresponding segment of the split model as a teacher for feature distillation~\cite{nasfly_2024}. This converts waiting time into architecture-search work that can proceed without the remaining model partitions.

\subsection{Limited Networking} % networking related challenges and solutions

Cross-device clients often communicate with the server over capacity-limited wide-area links. Typical wide-area bandwidths of 5--25 Mbit/s are far below the more than 10 Gbit/s available within data centres~\cite{peaches_2024}. The discussion below assumes a common bandwidth limit. Unequal connectivity is treated separately as client heterogeneity.

\subsubsection{Challenge 1.6: Search Messages Exceed Client Transfer Budgets}

Fixed-model FL exchanges the parameters of one architecture. NAS-in-FL may additionally exchange an over-parameterised supernet or architecture variables. A supernet contains weights for mutually exclusive operations and can therefore be much larger than any selected architecture, making the complete search state infeasible to transfer over a constrained client link~\cite{fedoras_2022,hapfnas_2025}.

\noindent\textbf{Solution 1.6a: Compress search updates.} Quantisation represents update entries with fewer bits, while sparsification transmits only selected entries. One differentiable supernet-based approach uses 8-bit stochastic quantisation and retains the largest 20\% of entries after clipping~\cite{fedregnas_2025}. These operations reduce the transferred weight and architecture updates without changing the search variables.

\noindent\textbf{Solution 1.6b: Encode architecture choices compactly.} When clients already hold a compatible master model or trained supernet, the server can identify a candidate using a selector or genotype instead of transmitting its parameters. Clients reconstruct the indicated subnet locally and return its evaluation or updated weights~\cite{rt_fed_evo_nas_2020,medical_image_segmentation_fednas_2024,intrusion_detection_fednas_2024}. This requires the server and clients to share the search-space definition and clients to retain the common weights.

\subsubsection{Challenge 1.7: Architecture Search Adds Repeated Synchronisation}

Differentiable NAS-in-FL updates architecture variables alongside model weights. Synchronising both in every round adds a communication channel that fixed-model FL does not require~\cite{fednas_2021,fedregnas_2025}. Even small architecture updates can consume substantial bandwidth when exchanged throughout a long search.

\noindent\textbf{Solution 1.7a: Synchronise architecture variables less frequently.} Architecture variables can be treated as slow variables that are updated and uploaded less frequently than model weights. One approach transmits them every \(H_{\alpha}\) rounds and stops doing so after the architecture has been fixed, while ordinary weight training continues~\cite{fedregnas_2025}.

\noindent\textbf{Solution 1.7b: Aggregate search updates hierarchically.} A nearby aggregator can combine the model and architecture updates of several clients before forwarding one result to the remote server. Frequent aggregation then uses local links, while fewer exchanges traverse the wide-area network~\cite{finch_2024}. In one non-IID CIFAR-10 evaluation, this reduced traffic at 80\% target accuracy from 3.29--4.68 GB to 1.64 GB~\cite{finch_2024}.

\noindent\textbf{Solution 1.7c: Parallelise architecture and weight optimisation.} Vertically partitioned differentiable strategies can formulate the architecture and weight updates independently within a round, allowing both to proceed in parallel rather than through alternating steps. This removes a sequential dependency and reduces the synchronised exchanges needed to update both variable sets~\cite{self-supervised_vertical_fednas_2023}.

\subsubsection{Challenge 1.8: Communication-Unaware Search Produces Expensive Federated Models}

An architecture selected only for predictive performance or local computation can still contain many parameters that FL must transmit in every subsequent training round. The search decision therefore determines the continuing load on client links, and a model that meets computation and memory budgets may still exceed the transfer budget~\cite{multi_objective_evo_fl_2020}.

\noindent\textbf{Solution 1.8a: Optimise for transmitted model size.} Communication cost can be included directly in the architecture objective. Evolutionary search can optimise global error together with the average number of weights uploaded by clients, favouring sparse architectures~\cite{multi_objective_evo_fl_2020}. Differentiable supernet-based search can similarly regularise the expected number of transmitted bytes~\cite{fedregnas_2025}. Both expose the trade-off between predictive performance and continuing communication cost before selecting an architecture.

\subsection{Energy Constraints}

\subsubsection{Challenge 1.9: Search Computation and Communication Deplete Client Energy}

Message size and round duration do not fully represent the burden on battery-powered clients. Repeated local optimisation consumes computation energy, while every search-state exchange consumes transmission energy. Fixed local epoch counts can create redundant communication in some rounds and unnecessary computation in others, increasing the energy required to reach a target performance~\cite{f2nas_2024}.

\noindent\textbf{Solution 1.9a: Adapt local epochs to an energy--convergence bound.} The number of local epochs can be selected in each round from an upper bound that relates loss reduction to client computation and communication costs. This balances additional local work against another cloud--client exchange instead of using one fixed value throughout the search~\cite{f2nas_2024}. The cited evaluation reports an energy reduction of nearly 50\%, although applying the mechanism requires estimates of client energy and latency.

\subsection{Client Unreliability}

\subsubsection{Challenge 1.10: Client Dropout Interrupts Architecture Search}

Cross-device clients can become unavailable because of changing network conditions, battery constraints or competing local workloads. A search procedure that requires the same cohort in every round can then stall or lose updates for the candidates or supernet paths assigned to disconnected clients~\cite{optimize_fednas_edge_2021}.

\noindent\textbf{Solution 1.10a: Design the search for partial participation.} The server can require only a random subset of clients to participate in each round and aggregate the search updates returned by that subset. The participating clients may change between rounds, allowing architecture search to continue without requiring every registered client to remain available~\cite{optimize_fednas_edge_2021}.

\subsubsection{Challenge 1.11: Clients Reach Architecture Decisions at Different Times}

In personalised differentiable search, clients can optimise architecture variables at different rates because their data distributions and local trajectories differ. A common stopping round can terminate some searches prematurely while allowing others to overfit, for example by increasingly favouring skip connections~\cite{collaborative_nas_for_pfl_2025}.

\noindent\textbf{Solution 1.11a: Stop architecture updates independently.} After a shared warm-up period, each client can stop updating its architecture variables when a local criterion detects increasing preference for skip connections. Other clients continue searching until they meet their own stopping criteria~\cite{collaborative_nas_for_pfl_2025}.
